\chapter{Estado del Arte}
\label{chap:estado_del_arte}

El desarrollo de sistemas de trading algorítmico ha experimentado una transformación radical en la última década, pasando de la ejecución de simples reglas estáticas a la integración de modelos complejos de Inteligencia Artificial (IA) \parencite{treleaven_algorithmic_2013}. En esta sección se analizan las tecnologías, paradigmas y algoritmos que conforman el estado actual de la predicción de mercados financieros y las arquitecturas de software de trading. El recorrido de este capítulo sigue la misma lógica del sistema desarrollado: primero se sitúa la evolución del trading algorítmico, después cómo se construyen las variables de entrada del módulo predictivo, a continuación se revisan las familias de modelos que permiten formular la predicción direccional y, finalmente, se analiza la arquitectura software necesaria para integrar cálculo analítico, persistencia y ejecución.

\section{Evolución del Trading Algorítmico}

El trading algorítmico consiste en el uso de programas informáticos para introducir órdenes de compra y venta en los mercados financieros con base en un conjunto de reglas predefinidas \parencite{chan_algorithmic_2013}. Históricamente, estos sistemas se basaban en el Análisis Técnico tradicional, utilizando indicadores matemáticos retrospectivos \parencite{murphy_technical_1999}. Sin embargo, la Hipótesis de los Mercados Eficientes \parencite{fama_efficient_1970} sugiere que los precios de los activos reflejan toda la información disponible, lo que hace que los métodos puramente lineales sean insuficientes en mercados de alta liquidez como el de las criptomonedas \parencite{katsiampa_volatility_2017}. Esto ha provocado la migración hacia el análisis cuantitativo y el modelado no lineal mediante aprendizaje automático \parencite{lopez_de_prado_advances_2018}.

Esta evolución explica que, antes de discutir algoritmos concretos, resulte necesario revisar cómo se construyen las variables con las que esos modelos operan. En un sistema como el desarrollado en este trabajo, la predicción no se plantea sobre velas OHLCV en bruto, sino sobre indicadores que representan la tendencia, momento y volatilidad.

\section{Datos y Modelos Predictivos}

A pesar del auge de la IA, los indicadores técnicos clásicos siguen siendo la base fundamental para la fase de extracción de características \parencite{lopez_de_prado_advances_2018}. La literatura actual demuestra que la combinación de indicadores de tendencia (como la Media Móvil Exponencial, EMA), de momento (como el Relative Strength Index, RSI) y de volatilidad (como el Average True Range, ATR) proporcionan a los modelos predictivos el contexto necesario para identificar regímenes de mercado \parencite[Caps. 9, 10 y 16]{murphy_technical_1999}. El preprocesamiento de estos datos en ventanas temporales escalonadas es vital para reducir el ruido del precio \parencite{chan_algorithmic_2013,lopez_de_prado_advances_2018}.

En este proyecto, el módulo predictivo se apoya precisamente en un conjunto de variables calculables y comparables. Una vez construido ese espacio de características, el problema pasa a ser decidir qué familias de modelos supervisados resultan más adecuadas para predecir la dirección de la siguiente vela. La predicción del movimiento direccional del precio (clasificación: sube, baja o lateral) es uno de los problemas más complejos en el ámbito del aprendizaje automático, debido a la baja relación señal-ruido de los mercados financieros \parencite{fama_efficient_1970}. Actualmente, la industria se apoya en las siguientes familias de algoritmos:

\begin{description}
    \item[\normalfont\textbf{Modelos de Ensamblado basados en Árboles.}] Los métodos de ensamblado dominan las competiciones de ciencia de datos tabulares financieras \parencite{lopez_de_prado_advances_2018}. \textbf{Random Forest} actúa como el estándar de la industria para establecer referencias robustas gracias a su resistencia al sobreajuste, lograda al combinar múltiples árboles de decisión independientes \parencite{breiman_random_2001}. Por su parte, los métodos de \textbf{Gradient Boosting} como XGBoost \parencite{chen_xgboost_2016} y LightGBM \parencite{ke_lightgbm_2017} destacan por su alta precisión en conjuntos de datos tabulares: XGBoost incorpora técnicas de regularización que reducen el sobreajuste y mejoran la selección de variables relevantes \parencite{chen_xgboost_2016}, mientras que LightGBM ofrece un entrenamiento más rápido y eficiente en memoria \parencite{ke_lightgbm_2017}, siendo el algoritmo preferido para procesar grandes volúmenes de datos históricos.

    \item[\normalfont\textbf{Máquinas de Vectores de Soporte (SVM).}] En su formulación matemática, las SVM transforman los datos de entrada a un espacio de mayor dimensión mediante funciones núcleo para encontrar un hiperplano de separación óptimo \parencite{cortes_support-vector_1995}. Aplicado al dominio financiero, aunque estas técnicas resultan computacionalmente costosas al procesar conjuntos de datos masivos, siguen siendo una herramienta valiosa para identificar cambios de tendencia macroeconómicos cuando se utilizan con variables muy refinadas \parencite{treleaven_algorithmic_2013,pedregosa_scikit_2011}.

    \item[\normalfont\textbf{Redes Neuronales y Aprendizaje Profundo.}] Las redes neuronales profundas han demostrado una capacidad excepcional para aprender representaciones no lineales complejas en datos financieros \parencite{goodfellow_deep_2016}. Cuando el problema se formula como clasificación direccional sobre un conjunto de indicadores técnicos precalculados, las arquitecturas densas totalmente conectadas resultan la elección natural: el modelo recibe los indicadores como entrada y predice la dirección de la siguiente vela \parencite{goodfellow_deep_2016}. La arquitectura típica combina capas de neuronas con funciones de activación no lineales, capas de regularización para reducir el sobreajuste y una capa de salida para clasificación multiclase. Su principal ventaja sobre los modelos de ensamblado radica en la capacidad de gestionar interacciones no lineales de alto orden entre indicadores \parencite{goodfellow_deep_2016}; sin embargo, su mayor tendencia al sobreajuste en conjuntos de datos reducidos y su coste computacional de entrenamiento los convierten en una herramienta complementaria a los métodos basados en árboles \parencite{lopez_de_prado_advances_2018}.
\end{description}

\section{Arquitecturas de Software para Sistemas de Trading}

Una vez fijados el tipo de variables y las familias de modelos predictivos, la cuestión deja de ser solo algorítmica y pasa a ser también de integración: cómo combinar entrenamiento, inferencia, ingesta de datos y ejecución transaccional sin mezclar responsabilidades. Ese es el contexto en el que la arquitectura software adquiere relevancia para este trabajo.

El ecosistema del desarrollo de bots de trading requiere una arquitectura distribuida que soporte alta concurrencia y baja latencia \parencite{treleaven_algorithmic_2013}. El estado del arte en ingeniería de software financiera divide estas soluciones en dos capas \parencite{richards_software_2022}:

\begin{enumerate}
    \item \textbf{Capa Transaccional y de Orquestación:} Generalmente desarrollada en lenguajes fuertemente tipados y compilados orientados a concurrencia, como Java con entornos como Spring Boot \parencite{pivotal_springboot_2024} o C++. Esta capa gestiona el enrutamiento de red, la persistencia masiva en bases de datos (inserciones masivas para series temporales) \parencite{kimball_data_2004} y la ingesta de datos en tiempo real \parencite{richards_software_2022}.

    \item \textbf{Capa Analítica (Inferencia y Entrenamiento):} Implementada casi en su totalidad en Python \parencite{hilpisch_python_2018} debido a su ecosistema superior en ciencia de datos.
\end{enumerate}

Dos de los problemas más comunes y críticos en este tipo de arquitecturas son la comunicación entre capas y la configuración de los modelos predictivos. Por un lado, la comunicación entre procesos (IPC) representa uno de los mayores desafíos técnicos \parencite{furuhashi_msgpack_2013}: las soluciones actuales optan por colas de mensajes, servicios web o comunicación directa mediante tuberías y entrada/salida estándar \parencite{treleaven_algorithmic_2013,newman_building_2015,kleppmann_designing_2017}. Por otro lado, ajustar los parámetros de los modelos predictivos de forma manual resulta inviable cuando el espacio de posibilidades es amplio; por ello, el estado del arte emplea métodos de optimización bayesiana que guían la búsqueda de forma eficiente \parencite{lopez_de_prado_advances_2018}, reduciendo el tiempo necesario sin sacrificar la calidad de la configuración encontrada \parencite{feurer_hyperparameter_2019}.

\section{Posicionamiento del Trabajo}

Sobre esta base, el posicionamiento del trabajo puede formularse con mayor precisión: no se trata solo de construir otro bot de trading, sino de proponer una plataforma que combine el estado del arte en variables técnicas, modelos predictivos y desacoplamiento arquitectónico dentro de un entorno docente y reproducible.

Para comprender el valor de la plataforma desarrollada en este proyecto, es necesario analizar las soluciones de trading algorítmico existentes en el mercado actual y sus limitaciones, especialmente desde una perspectiva académica y educativa \parencite{lusardi_economic_2014}.

En el ámbito comercial, herramientas líderes como \textit{MetaTrader} o \textit{TradingView} dominan el sector minorista. Sin embargo, presentan barreras significativas: operan frecuentemente como ``cajas negras'' donde el motor de ejecución es opaco, o requieren el uso de lenguajes de programación propietarios y de nicho (como \textit{MQL} o \textit{PineScript}). Estos lenguajes carecen de integración nativa con el ecosistema moderno de Inteligencia Artificial \parencite{goodfellow_deep_2016}, lo que obliga a los usuarios a depender de complejas pasarelas de red para conectar sus modelos predictivos \parencite{newman_building_2015}. Por otro lado, plataformas en la nube como \textit{3Commas} o \textit{Cryptohopper} facilitan la automatización, pero a costa de suscripciones mensuales y una nula capacidad de personalización algorítmica profunda \parencite{hasso_who_2019}.

En el espectro del software de código abierto, existen proyectos muy consolidados. Plataformas como \textit{QuantConnect} ofrecen una potencia institucional, pero su altísima complejidad técnica y su dependencia de infraestructuras en la nube las hacen poco accesibles para estudiantes que buscan un primer contacto con el trading cuantitativo \parencite{lusardi_economic_2014}. Por su parte, \textit{Freqtrade} \parencite{freqtrade_development_team_freqtrade_2023} es una herramienta excelente y accesible, pero su arquitectura es monolítica (puramente en Python) \parencite{freqtrade_team_freqtrade_2026}. Esto significa que la misma aplicación gestiona la red, la base de datos, la contabilidad y la IA, lo que no representa el estándar de la industria del desarrollo de software empresarial, donde estos componentes suelen estar desacoplados \parencite{richards_software_2022}.

Frente a estas alternativas, este Trabajo de Fin de Grado se posiciona y diferencia al proponer una arquitectura híbrida, modular y con un enfoque netamente educativo \parencite{chan_algorithmic_2013}. A diferencia de las opciones comerciales, el sistema es transparente, local y utiliza Python como lenguaje principal para la formulación de estrategias, abriendo la puerta de forma nativa al estado del arte del aprendizaje automático \parencite{lopez_de_prado_advances_2018}. A diferencia de las soluciones de código abierto monolíticas, este proyecto separa drásticamente las responsabilidades: un orquestador Java \parencite{pivotal_springboot_2024} garantiza la robustez, concurrencia e integridad transaccional típica de los sistemas bancarios \parencite{oracle_mysql_acid_2024} (gestionando el libro mayor y la operativa simulada \parencite{chan_algorithmic_2013}), mientras que subprocesos independientes de Python actúan únicamente como el motor matemático \parencite{hilpisch_python_2018}.

Esta clara separación no solo resuelve el problema técnico de la ingesta masiva de datos y su persistencia segura \parencite{kimball_data_2004}, sino que proporciona al estudiante una herramienta de experimentación inmejorable \parencite{lusardi_economic_2014}: expone una interfaz flexible capaz de ajustar los parámetros de los distintos modelos predictivos integrados, e incorpora además un módulo de búsqueda automática de la configuración óptima basado en optimización bayesiana \parencite{feurer_hyperparameter_2019}, todo ello en un entorno seguro y replicable \parencite{lopez_de_prado_advances_2018}.